\subsection{Radicals, Hyperbolic Spaces, and Witt Theory}

Throughout this subsection, metric vector spaces are finite-dimensional over a fixed arbitrary field $\mathbb{F}$.

\begin{definition} \label{def:metric-orthogonal-direct-sum} \hlt{Orthogonal Direct Sum}\\
For $\mathcal{S},\mathcal{T}\leq\mathcal{V}$, the notation $\mathcal{V}=\mathcal{S}\odot\mathcal{T}$ means $\mathcal{V}=\mathcal{S}\oplus\mathcal{T}$ and $\mathcal{S}\perp\mathcal{T}$; then $\mathcal{V}$ is the \hlt{orthogonal direct sum} of $\mathcal{S}$ and $\mathcal{T}$.
\end{definition}

\begin{theorem} \label{thm:metric-orthogonal-complement-dimension} \hlt{Dimensions of Orthogonal Complements}\\
Let $\mathcal{S}\leq\mathcal{V}$.
\begin{enumerate}[label=(\roman*)]
\setlength{\itemsep}{0pt}
\item If $\operatorname{rad}(\mathcal{V})=\{\mathbf{0}\}$ or $\operatorname{rad}(\mathcal{S})=\{\mathbf{0}\}$, then $\dim(\mathcal{S})+\dim(\mathcal{S}^{\perp})=\dim(\mathcal{V})$.
\item If $\operatorname{rad}(\mathcal{S})=\{\mathbf{0}\}$, then $\mathcal{V}=\mathcal{S}\odot\mathcal{S}^{\perp}$.
\item If $\operatorname{rad}(\mathcal{V})=\{\mathbf{0}\}$, then $\mathcal{S}^{\perp\perp}=\mathcal{S}$.
\item If $\operatorname{rad}(\mathcal{V})=\{\mathbf{0}\}$, then $\operatorname{rad}(\mathcal{S})=\{\mathbf{0}\}$ $\Longleftrightarrow$ $\operatorname{rad}(\mathcal{S}^{\perp})=\{\mathbf{0}\}$.
\end{enumerate}
\begin{proof}
\begin{enumerate}[label=(\roman*)]
\setlength{\itemsep}{0pt}
\item[$(i)$] Define $\rho:\mathcal{V}\to\mathcal{S}'$ by $\rho(\mathbf{x})(\mathbf{s})=\innerprod{\mathbf{s},\mathbf{x}}$. Bilinearity implies $\rho\in\LL(\mathcal{V},\mathcal{S}')$, and $\ker(\rho)=\mathcal{S}^{\perp}$.\\
Fix $f\in\mathcal{S}'$. If $\operatorname{rad}(\mathcal{V})=\{\mathbf{0}\}$, choose $\mathcal{V}=\mathcal{S}\oplus\mathcal{C}$ by Proposition \ref{prop:complements-exist-finite-dimensions} and define $\widetilde f(\mathbf{s}+\mathbf{c})=f(\mathbf{s})$.\\
By Theorem \ref{thm:riesz-map-for-nonsingular-metric-spaces}, fix $\mathbf{x}\in\mathcal{V}$ satisfying $(\forall\,\mathbf{s}\in\mathcal{S})\,f(\mathbf{s})=\widetilde f(\mathbf{s})=\innerprod{\mathbf{s},\mathbf{x}}$. If $\operatorname{rad}(\mathcal{S})=\{\mathbf{0}\}$, fix such an $\mathbf{x}\in\mathcal{S}$ by Theorem \ref{thm:riesz-map-for-nonsingular-metric-spaces} for $\mathcal{S}$.\\
Thus $\rho$ is surjective in either case. By Theorem \ref{thm:rank-nullity} and Theorem \ref{thm:dual-basis-and-dimension}(ii), $\dim(\mathcal{V})=\dim(\ker\rho)+\dim(\mathcal{S}')=\dim(\mathcal{S}^{\perp})+\dim(\mathcal{S})$.
\item[$(ii)$] If $\operatorname{rad}(\mathcal{S})=\{\mathbf{0}\}$, then $\mathcal{S}\cap\mathcal{S}^{\perp}=\{\mathbf{0}\}$. This trivial intersection and (i) imply $\mathcal{V}=\mathcal{S}\oplus\mathcal{S}^{\perp}$.\\
Also $\mathcal{S}\perp\mathcal{S}^{\perp}$ holds by the definition of $\mathcal{S}^{\perp}$. Hence $\mathcal{V}=\mathcal{S}\odot\mathcal{S}^{\perp}$.
\item[$(iii)$] Suppose that $\operatorname{rad}(\mathcal{V})=\{\mathbf{0}\}$. The inclusion $\mathcal{S}\subseteq\mathcal{S}^{\perp\perp}$ is immediate.\\
By (i) for $\mathcal{S}$ and $\mathcal{S}^{\perp}$, $\dim(\mathcal{S}^{\perp\perp})=\dim(\mathcal{S})$, so $\mathcal{S}^{\perp\perp}=\mathcal{S}$ by Proposition \ref{prop:dimension-bounds-and-basis-criteria}(ii).
\item[$(iv)$] ($\Rightarrow$) If $\operatorname{rad}(\mathcal{S})=\{\mathbf{0}\}$, then $\operatorname{rad}(\mathcal{S}^{\perp})=\mathcal{S}^{\perp}\cap\mathcal{S}^{\perp\perp}=\mathcal{S}^{\perp}\cap\mathcal{S}=\{\mathbf{0}\}$ by (iii).
\item[$(iv)$] ($\Leftarrow$) If $\operatorname{rad}(\mathcal{S}^{\perp})=\{\mathbf{0}\}$, then $\operatorname{rad}(\mathcal{S})=\mathcal{S}\cap\mathcal{S}^{\perp}=\mathcal{S}^{\perp\perp}\cap\mathcal{S}^{\perp}=\{\mathbf{0}\}$ by (iii).
\end{enumerate}
\hlt{Structural Commentary}: A restriction pairing gives the dimension formula; dimension counting gives the rest.\\
\hlt{Key Insight}: Nonsingularity turns orthogonal complementation into a dimension-reversing involution.
\end{proof}
\end{theorem}

\begin{theorem} \label{thm:radical-splitting-metric-space} \hlt{Splitting off the Radical}\\
Let $R=\operatorname{rad}(\mathcal{V})$. Then $\exists\,\mathcal{S}\leq\mathcal{V}$ such that $\mathcal{V}=R\odot\mathcal{S}$, $\operatorname{rad}(\mathcal{S})=\{\mathbf{0}\}$, and $R$ is totally singular.
\begin{proof}
By Proposition \ref{prop:complements-exist-finite-dimensions}, choose $\mathcal{S}$ satisfying $\mathcal{V}=R\oplus\mathcal{S}$. Since $R\perp\mathcal{V}$, $\mathcal{V}=R\odot\mathcal{S}$ and the restricted form on $R$ vanishes.\\
Thus $R$ is totally singular. If $\mathbf{s}\in\operatorname{rad}(\mathcal{S})$, then $\mathbf{s}\perp\mathcal{S}$ and $\mathbf{s}\perp R$, so $\mathbf{s}\in R\cap\mathcal{S}=\{\mathbf{0}\}$.\\
\hlt{Structural Commentary}: A vector-space complement isolates the radical and inherits a nonsingular restricted form.\\
\hlt{Key Insight}: Every singular direction is confined to an orthogonal zero-form summand.
\end{proof}
\end{theorem}

\begin{proposition} \label{prop:metric-quotients-by-radical-subspaces} \hlt{Metric Quotients}\\
Let $\mathcal{U}\leq\mathcal{V}$ and consider the rule $\innerprod[\mathcal{V}/\mathcal{U}]{\mathbf{x}+\mathcal{U},\mathbf{y}+\mathcal{U}}\defeq\innerprod{\mathbf{x},\mathbf{y}}$.
\begin{enumerate}[label=(\roman*)]
\setlength{\itemsep}{0pt}
\item The rule is well-defined $\Longleftrightarrow$ $\mathcal{U}\leq\operatorname{rad}(\mathcal{V})$.
\item If $\mathcal{U}\leq\operatorname{rad}(\mathcal{V})$, then $\operatorname{rad}(\mathcal{V}/\mathcal{U})=\operatorname{rad}(\mathcal{V})/\mathcal{U}$.
\item The metric space $\mathcal{V}/\operatorname{rad}(\mathcal{V})$ is nonsingular.
\end{enumerate}
\begin{proof}
\begin{enumerate}[label=(\roman*)]
\setlength{\itemsep}{0pt}
\item[$(i)$] ($\Rightarrow$) Suppose that the rule is well-defined. Then $(\forall\,\mathbf{u}\in\mathcal{U})(\forall\,\mathbf{y}\in\mathcal{V})\,0=\innerprod[\mathcal{V}/\mathcal{U}]{\mathbf{u}+\mathcal{U},\mathbf{y}+\mathcal{U}}=\innerprod{\mathbf{u},\mathbf{y}}$, so $\mathcal{U}\leq\operatorname{rad}(\mathcal{V})$.
\item[$(i)$] ($\Leftarrow$) Suppose that $\mathcal{U}\leq\operatorname{rad}(\mathcal{V})$. By Proposition \ref{prop:metric-radical-equals-bilinear-radicals}(i), $(\forall\,\mathbf{u},\mathbf{v}\in\mathcal{U})(\forall\,\mathbf{x},\mathbf{y}\in\mathcal{V})\,\innerprod{\mathbf{x}+\mathbf{u},\mathbf{y}+\mathbf{v}}=\innerprod{\mathbf{x},\mathbf{y}}$. QED.
\item[$(ii)$] Suppose that $\mathcal{U}\leq\operatorname{rad}(\mathcal{V})$. Then $\mathbf{x}+\mathcal{U}\in\operatorname{rad}(\mathcal{V}/\mathcal{U})$ $\Longleftrightarrow$ $(\forall\,\mathbf{y}\in\mathcal{V})\,\innerprod{\mathbf{x},\mathbf{y}}=0$ $\Longleftrightarrow$ $\mathbf{x}\in\operatorname{rad}(\mathcal{V})$. Thus $\operatorname{rad}(\mathcal{V}/\mathcal{U})=\operatorname{rad}(\mathcal{V})/\mathcal{U}$.
\item[$(iii)$] Set $\mathcal{U}=\operatorname{rad}(\mathcal{V})$ in (ii). Then $\operatorname{rad}(\mathcal{V}/\mathcal{U})=\operatorname{rad}(\mathcal{V})/\operatorname{rad}(\mathcal{V})=\{\mathbf{0}\}$.
\end{enumerate}
\hlt{Structural Commentary}: Item (i) tests representative independence; items (ii)--(iii) compute and remove the residual radical.\\
\hlt{Key Insight}: Exactly radical directions may be collapsed without changing pairings.
\end{proof}
\end{proposition}

\begin{proposition} \label{prop:orthogonal-sums-radicals-and-regular-quotients} \hlt{Radicals and Orthogonal Direct Sums}\\
Suppose that $\mathcal{V}=\mathcal{S}\odot\mathcal{T}$.
\begin{enumerate}[label=(\roman*)]
\setlength{\itemsep}{0pt}
\item $\operatorname{rad}(\mathcal{V})=\operatorname{rad}(\mathcal{S})\odot\operatorname{rad}(\mathcal{T})$.
\item $\mathcal{V}/\operatorname{rad}(\mathcal{V})\cong(\mathcal{S}/\operatorname{rad}(\mathcal{S}))\odot(\mathcal{T}/\operatorname{rad}(\mathcal{T}))$.
\item $\mathcal{V}$ is nonsingular $\Longleftrightarrow$ $\mathcal{S}$ and $\mathcal{T}$ are nonsingular.
\end{enumerate}
\begin{proof}
\begin{enumerate}[label=(\roman*)]
\setlength{\itemsep}{0pt}
\item[$(i)$] Fix $\mathbf{v}=\mathbf{s}+\mathbf{t}$ with $\mathbf{s}\in\mathcal{S}$ and $\mathbf{t}\in\mathcal{T}$. Orthogonality implies $\mathbf{v}\in\operatorname{rad}(\mathcal{V})$ $\Longleftrightarrow$ $\mathbf{s}\in\operatorname{rad}(\mathcal{S})$ and $\mathbf{t}\in\operatorname{rad}(\mathcal{T})$.\\
Uniqueness of the decomposition proves the radical identity.
\item[$(ii)$] Define $\Phi((\mathbf{s}+\mathbf{t})+\operatorname{rad}(\mathcal{V}))=(\mathbf{s}+\operatorname{rad}(\mathcal{S}),\mathbf{t}+\operatorname{rad}(\mathcal{T}))$. By (i), $\Phi$ well-defined, has inverse $\Psi(\mathbf{s}+\operatorname{rad}(\mathcal{S}),\mathbf{t}+\operatorname{rad}(\mathcal{T}))=(\mathbf{s}+\mathbf{t})+\operatorname{rad}(\mathcal{V})$. Hence $\Phi$ is linear bijection. By orthogonality, $(\forall\,\mathbf{s}_1,\mathbf{s}_2\in\mathcal{S})(\forall\,\mathbf{t}_1,\mathbf{t}_2\in\mathcal{T})\,\innerprod{\mathbf{s}_1+\mathbf{t}_1,\mathbf{s}_2+\mathbf{t}_2}=\innerprod{\mathbf{s}_1,\mathbf{s}_2}+\innerprod{\mathbf{t}_1,\mathbf{t}_2}$; $\Phi$ is isometry.
\item[$(iii)$] ($\Rightarrow$) If $\mathcal{V}$ is nonsingular, then $\operatorname{rad}(\mathcal{S})\odot\operatorname{rad}(\mathcal{T})=\{\mathbf{0}\}$ by (i), so $\mathcal{S}$ and $\mathcal{T}$ are nonsingular.
\item[$(iii)$] ($\Leftarrow$) If $\mathcal{S}$ and $\mathcal{T}$ are nonsingular, then $\operatorname{rad}(\mathcal{V})=\{\mathbf{0}\}$ by (i).
\end{enumerate}
\hlt{Structural Commentary}: Item (i) splits the radical; item (ii) descends the decomposition; item (iii) detects nonsingularity.\\
\hlt{Key Insight}: Orthogonality separates both the form and its invisible directions.
\end{proof}
\end{proposition}

\begin{corollary} \label{cor:metric-spaces-determined-by-regular-quotient-and-dimension} \hlt{Recovering a Singular Space from Its Regular Quotient}\\
If $\dim\mathcal{V}=\dim\mathcal{W}$ and $\mathcal{V}/\operatorname{rad}(\mathcal{V}) \cong \mathcal{W}/\operatorname{rad}(\mathcal{W})$, then $\mathcal{V}\cong\mathcal{W}$ as metric vector spaces.
\begin{proof}
\begin{enumerate}[label=(\roman*)]
\setlength{\itemsep}{0pt}
\item \hlt{Split}: By Theorem \ref{thm:radical-splitting-metric-space}, choose $\mathcal{V}=\operatorname{rad}(\mathcal{V})\odot\mathcal{S}$ and $\mathcal{W}=\operatorname{rad}(\mathcal{W})\odot\mathcal{T}$ with $\mathcal{S},\mathcal{T}$ nonsingular.
\item \hlt{Regular Parts}: By Proposition \ref{prop:orthogonal-sums-radicals-and-regular-quotients}(ii), $\mathcal{V}/\operatorname{rad}(\mathcal{V})\cong\mathcal{S}$ and $\mathcal{W}/\operatorname{rad}(\mathcal{W})\cong\mathcal{T}$. Hence $\mathcal{S}\cong\mathcal{T}$ by the assumed quotient isometry.
\item \hlt{Radicals}: The dimension hypotheses imply $\dim\operatorname{rad}(\mathcal{V})=\dim\operatorname{rad}(\mathcal{W})$.\\
Fix a vector-space isomorphism between these radicals; it is an isometry because both restricted forms vanish.
\item \hlt{Assembly}: The orthogonal direct sum of the radical isometry and an isometry $\mathcal{S}\to\mathcal{T}$ is an isometry $\mathcal{V}\to\mathcal{W}$.
\end{enumerate}
\hlt{Structural Commentary}: Split both radicals, identify the regular parts, match dimensions, and reassemble orthogonally.\\
\hlt{Key Insight}: A radical carries only its dimension once its form vanishes.
\end{proof}
\end{corollary}

\begin{definition} \label{def:hyperbolic-pairs-and-spaces} \hlt{Hyperbolic Pairs and Hyperbolic Spaces}
\begin{enumerate}[label=(\roman*)]
\setlength{\itemsep}{0pt}
\item A \hlt{hyperbolic pair} is $(\mathbf{u},\mathbf{v})$ with $\innerprod{\mathbf{u},\mathbf{u}}=\innerprod{\mathbf{v},\mathbf{v}}=0$ and $\innerprod{\mathbf{u},\mathbf{v}}=1$.
\item The plane $\operatorname{span}(\mathbf{u},\mathbf{v})$ of a hyperbolic pair is a \hlt{hyperbolic plane}.
\item A \hlt{hyperbolic space} is an orthogonal direct sum of hyperbolic planes.
\item A basis $(\mathbf{u}_1,\mathbf{v}_1,\ldots,\mathbf{u}_k,\mathbf{v}_k)$ is a \hlt{hyperbolic basis} if each $(\mathbf{u}_i,\mathbf{v}_i)$ is hyperbolic and the planes $\operatorname{span}(\mathbf{u}_i,\mathbf{v}_i)$ are mutually orthogonal.\\
In the symplectic case it is a \hlt{symplectic basis}.
\end{enumerate}
\end{definition}

\begin{theorem} \label{thm:isotropic-vector-sits-in-hyperbolic-plane} \hlt{Capturing Isotropic Vectors}\\
Let $\mathcal{V}$ be nonsingular. If $\mathcal{V}$ is orthogonal, assume $\operatorname{char}(\mathbb{F})\neq2$.\\
Then $(\forall\,\mathbf{u}\in\mathcal{V}\setminus\{\mathbf{0}\})(\innerprod{\mathbf{u},\mathbf{u}}=0\Longrightarrow(\exists\,\mathbf{v}\in\mathcal{V})\,\innerprod{\mathbf{u},\mathbf{v}}=1\wedge\innerprod{\mathbf{v},\mathbf{v}}=0)$.
\begin{proof}
Fix $\mathbf{u}\neq\mathbf{0}$ with $\innerprod{\mathbf{u},\mathbf{u}}=0$. Since $\mathbf{u}\notin\operatorname{rad}(\mathcal{V})$, fix $\mathbf{x}\in\mathcal{V}$ with $c\defeq\innerprod{\mathbf{u},\mathbf{x}}\neq0$.\\
In alternating case, set $\mathbf{v}=c^{-1}\mathbf{x}$; then $\innerprod{\mathbf{u},\mathbf{v}}=1$ and $\innerprod{\mathbf{v},\mathbf{v}}=0$. In symmetric case, set $\mathbf{v}=c^{-1}\mathbf{x}-\frac{\innerprod{\mathbf{x},\mathbf{x}}}{2c^2}\mathbf{u}$.\\
Symmetry and $\innerprod{\mathbf{u},\mathbf{u}}=0$ imply $\innerprod{\mathbf{u},\mathbf{v}}=1$ and $\innerprod{\mathbf{v},\mathbf{v}}=c^{-2}\innerprod{\mathbf{x},\mathbf{x}}-c^{-2}\innerprod{\mathbf{x},\mathbf{x}}=0$. Thus $(\mathbf{u},\mathbf{v})$ is hyperbolic.\\
\hlt{Structural Commentary}: Choose a nonorthogonal partner and remove its self-pairing in the symmetric case.\\
\hlt{Key Insight}: One scalar correction converts a partner into an isotropic partner.
\end{proof}
\end{theorem}

\begin{theorem} \label{thm:nonsingular-completions-hyperbolic-extensions} \hlt{Nonsingular Completions}\\
Let $\mathcal{U}$ be a subspace of a nonsingular metric vector space $\mathcal{V}$, and assume $\operatorname{char}(\mathbb{F})\neq2$ in the orthogonal case.\\
Write $\mathcal{U}=\operatorname{span}(\mathbf{u}_1,\ldots,\mathbf{u}_r)\odot\mathcal{W}$, where $(\mathbf{u}_1,\ldots,\mathbf{u}_r)$ is a basis of $\operatorname{rad}(\mathcal{U})$ and $\mathcal{W}$ is nonsingular.
\begin{enumerate}[label=(\roman*)]
\setlength{\itemsep}{0pt}
\item $(\exists\,\mathbf{v}_1,\ldots,\mathbf{v}_r\in\mathcal{V})$ such that $\mathcal{H}_r=\operatorname{span}(\mathbf{u}_1,\mathbf{v}_1,\ldots,\mathbf{u}_r,\mathbf{v}_r)$ is hyperbolic and $\widehat{\mathcal{U}}=\mathcal{H}_r\odot\mathcal{W}$ contains $\mathcal{U}$ with $\dim\widehat{\mathcal{U}}=\dim\mathcal{U}+r$.
\item If $\mathcal{X}\supseteq\mathcal{U}$ is nonsingular, then $\dim\mathcal{X}\geq\dim\mathcal{U}+r$; equality is attained by $\widehat{\mathcal{U}}$, and all equality-case extensions are isometric.
\item If $\sigma:\mathcal{U}\to\mathcal{T}$ is an isometry between subspaces of nonsingular metric spaces of the same kind, equality-case extensions can be chosen so that $\sigma$ extends to an isometry between them.
\end{enumerate}
\begin{proof}
\begin{enumerate}[label=(\roman*)]
\setlength{\itemsep}{0pt}
\item[$(i)$] By Theorem \ref{thm:metric-orthogonal-complement-dimension}(ii), $\mathcal{V}=\mathcal{W}\odot\mathcal{W}^{\perp}$. Thus $(\forall\,i\in\{1,\ldots,r\})\,\mathbf{u}_i\in\mathcal{W}^{\perp}$. Set $\mathcal{H}_0=\{\mathbf{0}\}$ and proceed inductively.\\
After constructing $\mathcal{H}_{j-1}$, set $\mathcal{S}_j=\operatorname{span}(\mathbf{u}_{j+1},\ldots,\mathbf{u}_r)\odot\mathcal{W}\odot\mathcal{H}_{j-1}$. Pairing a relation with $\mathbf{u}_1,\ldots,\mathbf{u}_{j-1}$ eliminates its $\mathbf{v}_i$-coefficients, and independence of the radical basis modulo $\mathcal{W}$ then proves $\mathbf{u}_j\notin\mathcal{S}_j$.\\
By Theorem \ref{thm:metric-orthogonal-complement-dimension}(iii), $\mathcal{S}_j=\mathcal{S}_j^{\perp\perp}$, so fix $\mathbf{x}\in\mathcal{S}_j^{\perp}$ with $c\defeq\innerprod{\mathbf{u}_j,\mathbf{x}}\neq0$.\\
Set $\mathbf{v}_j=c^{-1}\mathbf{x}$ in the alternating case and $\mathbf{v}_j=c^{-1}\mathbf{x}-\frac{\innerprod{\mathbf{x},\mathbf{x}}}{2c^2}\mathbf{u}_j$ in the orthogonal case.\\
Then $\mathbf{v}_j\in\mathcal{S}_j^{\perp}$, $\innerprod{\mathbf{u}_j,\mathbf{v}_j}=1$, and $\innerprod{\mathbf{v}_j,\mathbf{v}_j}=0$ by direct calculation. Put $\mathcal{H}_j=\mathcal{H}_{j-1}\odot\operatorname{span}(\mathbf{u}_j,\mathbf{v}_j)$.\\
After $r$ steps, $\mathcal{H}_r$ is hyperbolic and $\widehat{\mathcal{U}}=\mathcal{H}_r\odot\mathcal{W}$ has the stated containment and dimension.
\item[$(ii)$] Let $\mathcal{X}\supseteq\mathcal{U}$ be nonsingular and choose $\mathcal{X}=\mathcal{U}\oplus\mathcal{Y}$. The pairing $\operatorname{rad}(\mathcal{U})\times\mathcal{Y}\to\mathbb{F}$ has zero left radical: a vector in that radical pairs trivially with both $\mathcal{U}$ and $\mathcal{Y}$ and hence lies in $\operatorname{rad}(\mathcal{X})$.\\
By Proposition \ref{prop:bilinear-radicals-matrix-and-dual-maps}(iii), the pairing matrix has rank $r$, so $\dim\mathcal{Y}\geq r$ and $\dim\mathcal{X}\geq\dim\mathcal{U}+r$. Item (i) attains equality.\\
If equality holds, the construction inside $\mathcal{X}$ produces a subspace of equal dimension, hence all of $\mathcal{X}$.\\
Any two equality-case extensions are isometric because both are $\mathcal{W}\odot$ a hyperbolic $r$-space.
\item[$(iii)$] Put $\mathbf{u}'_i=\sigma\mathbf{u}_i$ and $\mathcal{W}'=\sigma\mathcal{W}$. By (i) for $\mathcal{T}=\operatorname{span}(\mathbf{u}'_1,\ldots,\mathbf{u}'_r)\odot\mathcal{W}'$, choose partners $\mathbf{v}'_1,\ldots,\mathbf{v}'_r$.\\
Extend $\sigma$ by $(\forall\,i\in\{1,\ldots,r\})\,\mathbf{v}_i\mapsto\mathbf{v}'_i$. The resulting basis-to-basis linear map preserves the form on the hyperbolic and nonsingular summands, so it is an isometry between the equality-case extensions.
\end{enumerate}
\hlt{Structural Commentary}: Item (i) pairs radical directions; item (ii) proves minimality; item (iii) transports the construction.\\
\hlt{Key Insight}: Each radical dimension requires exactly one hyperbolic partner.
\end{proof}
\end{theorem}

\begin{definition} \label{def:witt-extension-cancellation-properties} \hlt{Witt Extension and Cancellation Properties}\\
Let $\mathcal{V}\cong\mathcal{V}'$ be nonsingular metric vector spaces of the same kind, either symplectic or orthogonal with $\operatorname{char}(\mathbb{F})\neq2$. Write $\tau:\mathcal{A}\xrightarrow{\sim}\mathcal{B}$ when $\tau$ is an isometry. The pair $(\mathcal{V},\mathcal{V}')$ has:
\begin{enumerate}[label=(\roman*)]
\setlength{\itemsep}{0pt}
\item the \hlt{Witt extension property} if $(\forall\,\mathcal{S}\leq\mathcal{V})(\forall\,\mathcal{T}\leq\mathcal{V}')(\forall\,\tau:\mathcal{S}\xrightarrow{\sim}\mathcal{T})(\exists\,\overline{\tau}:\mathcal{V}\xrightarrow{\sim}\mathcal{V}')\,\overline{\tau}|_{\mathcal{S}}=\tau$;
\item the \hlt{Witt cancellation property} if $(\forall\,\mathcal{S}\leq\mathcal{V})(\forall\,\mathcal{T}\leq\mathcal{V}')$, $\operatorname{rad}(\mathcal{S})=\operatorname{rad}(\mathcal{T})=\{\mathbf{0}\}$ and $\mathcal{S}\cong\mathcal{T}$ imply $\mathcal{S}^{\perp}\cong\mathcal{T}^{\perp}$.
\end{enumerate}
\end{definition}

\begin{proposition} \label{prop:witt-extension-cancellation-bridge} \hlt{Witt Extension--Cancellation Bridge}\\
Let $\mathcal{V}\cong\mathcal{V}'$ be nonsingular metric vector spaces of the same kind, either symplectic or orthogonal with $\operatorname{char}(\mathbb{F})\neq2$. Then the pair has the Witt extension property $\Longleftrightarrow$ it has the Witt cancellation property.
\begin{proof}
\begin{enumerate}[label=\roman*.]
\setlength{\itemsep}{0pt}
\item[$(\Rightarrow)$] Let $\mathcal{S}\leq\mathcal{V}$ and $\mathcal{T}\leq\mathcal{V}'$ be nonsingular, and fix an isometry $\sigma:\mathcal{S}\to\mathcal{T}$. By Theorem \ref{thm:metric-orthogonal-complement-dimension}(ii), $\mathcal{V}=\mathcal{S}\odot\mathcal{S}^{\perp}$ and $\mathcal{V}'=\mathcal{T}\odot\mathcal{T}^{\perp}$.\\
By Witt extension property, $\sigma$ extends to isometry $\widehat{\sigma}:\mathcal{V}\to\mathcal{V}'$. Form preservation, bijectivity imply $\widehat{\sigma}(\mathcal{S}^{\perp})=\mathcal{T}^{\perp}$, so $\mathcal{S}^{\perp}\cong\mathcal{T}^{\perp}$.
\item[$(\Leftarrow)$] Fix an isometry $\sigma:\mathcal{S}\to\mathcal{T}$ between arbitrary subspaces.\\
By Theorem \ref{thm:nonsingular-completions-hyperbolic-extensions}(iii), choose equality-case nonsingular extensions $\widehat{\mathcal{S}}\supseteq\mathcal{S}$ and $\widehat{\mathcal{T}}\supseteq\mathcal{T}$ and an isometry $\widehat{\sigma}:\widehat{\mathcal{S}}\to\widehat{\mathcal{T}}$ extending $\sigma$.\\
By Theorem \ref{thm:metric-orthogonal-complement-dimension}(ii), $\mathcal{V}=\widehat{\mathcal{S}}\odot\widehat{\mathcal{S}}^{\perp}$ and $\mathcal{V}'=\widehat{\mathcal{T}}\odot\widehat{\mathcal{T}}^{\perp}$. By the Witt cancellation property, fix an isometry $\eta:\widehat{\mathcal{S}}^{\perp}\to\widehat{\mathcal{T}}^{\perp}$.\\
Then $\widehat{\sigma}\odot\eta:\mathcal{V}\to\mathcal{V}'$ is an isometry extending $\sigma$.
\end{enumerate}
\hlt{Structural Commentary}: Extension restricts to complements; cancellation extends maps after nonsingular completion.\\
\hlt{Key Insight}: Completion reduces arbitrary subspaces to an orthogonal complement problem.
\end{proof}
\end{proposition}

\begin{theorem} \label{thm:witt-extension-theorem-metric-spaces} \hlt{Witt Extension Theorem}\\
Let $\mathcal{V}\cong\mathcal{W}$ be nonsingular metric vector spaces of the same kind, either symplectic or orthogonal with $\operatorname{char}(\mathbb{F})\neq2$.\\
If $\mathcal{S}\leq\mathcal{V}$, $\mathcal{T}\leq\mathcal{W}$, and $\sigma:\mathcal{S}\xrightarrow{\sim}\mathcal{T}$, then $(\exists\,\widehat{\sigma}:\mathcal{V}\xrightarrow{\sim}\mathcal{W})\,\widehat{\sigma}|_{\mathcal{S}}=\sigma$.
\begin{proof}
\begin{enumerate}[label=(\roman*)]
\setlength{\itemsep}{0pt}
\item \hlt{Symplectic Preparation}: Every two nonsingular symplectic spaces of equal dimension are isometric.\\
Indeed, a nonzero vector in such a space is isotropic and belongs to a hyperbolic plane by Theorem \ref{thm:isotropic-vector-sits-in-hyperbolic-plane}.\\
The plane splits off by Theorem \ref{thm:metric-orthogonal-complement-dimension}(ii), and its orthogonal complement is nonsingular by Theorem \ref{thm:metric-orthogonal-complement-dimension}(iv).\\
Induction on dimension produces a hyperbolic basis, so equal-dimensional spaces have isometric hyperbolic bases.
\item \hlt{Orthogonal Line Extension}: Let $\mathcal{V},\mathcal{W}$ be orthogonal, fix an isometry $\theta:\mathcal{V}\to\mathcal{W}$, and suppose $\innerprod{\mathbf{s},\mathbf{s}}=\innerprod{\mathbf{t},\mathbf{t}}\neq0$.\\
Put $\mathbf{x}=\theta\mathbf{s}$ and $\mathbf{y}=\mathbf{t}$. For nonisotropic $\mathbf{a}\in\mathcal{W}$, define $r_{\mathbf{a}}(\mathbf{z})=\mathbf{z}-\frac{2\innerprod{\mathbf{z},\mathbf{a}}}{\innerprod{\mathbf{a},\mathbf{a}}}\mathbf{a}$.\\
Direct expansion proves that $r_{\mathbf{a}}$ is an isometry and $r_{\mathbf{a}}\mathbf{a}=-\mathbf{a}$. If $\innerprod{\mathbf{x}-\mathbf{y},\mathbf{x}-\mathbf{y}}\neq0$, then $r_{\mathbf{x}-\mathbf{y}}\mathbf{x}=\mathbf{y}$.\\
Otherwise $\innerprod{\mathbf{x}+\mathbf{y},\mathbf{x}+\mathbf{y}}=4\innerprod{\mathbf{y},\mathbf{y}}\neq0$ and $r_{\mathbf{y}}r_{\mathbf{x}+\mathbf{y}}\mathbf{x}=\mathbf{y}$. Thus $r_{\mathbf{x}-\mathbf{y}}\theta$ or $r_{\mathbf{y}}r_{\mathbf{x}+\mathbf{y}}\theta$ extends $\mathbf{s}\mapsto\mathbf{t}$.
\item \hlt{Nonsingular Subspaces}: Assume that $\mathcal{S},\mathcal{T}$ are nonsingular and induct on $\dim\mathcal{S}$. The zero-dimensional case follows from $\mathcal{V}\cong\mathcal{W}$.\\
In the symplectic case, choose a hyperbolic plane $\mathcal{A}\leq\mathcal{S}$ by Theorem \ref{thm:isotropic-vector-sits-in-hyperbolic-plane}.\\
In the orthogonal case, fix $\mathbf{s}\in\mathcal{S}$ with $\innerprod{\mathbf{s},\mathbf{s}}\neq0$ by Proposition \ref{prop:form-type-relations-by-characteristic}(ii) and nonsingularity, and set $\mathcal{A}=\operatorname{span}(\mathbf{s})$. Put $\mathcal{B}=\sigma\mathcal{A}$.\\
In the symplectic case, fix an isometry $\beta:\mathcal{A}^{\perp}\to\mathcal{B}^{\perp}$ by (i) and set $\alpha=(\sigma|_{\mathcal{A}})\odot\beta$.\\
In the orthogonal case, fix an ambient isometry $\alpha$ extending $\sigma|_{\mathcal{A}}$ by (ii). Set $\mathcal{S}_0=\mathcal{S}\cap\mathcal{A}^{\perp}$ and $\mathcal{T}_0=\mathcal{T}\cap\mathcal{B}^{\perp}$.\\
By Theorem \ref{thm:metric-orthogonal-complement-dimension}(ii), $\mathcal{S}=\mathcal{A}\odot\mathcal{S}_0$, $\mathcal{T}=\mathcal{B}\odot\mathcal{T}_0$, $\mathcal{V}=\mathcal{A}\odot\mathcal{A}^{\perp}$, and $\mathcal{W}=\mathcal{B}\odot\mathcal{B}^{\perp}$.\\
The restriction $\alpha|_{\mathcal{A}^{\perp}}$ identifies the ambient complements, while $\dim\mathcal{S}_0<\dim\mathcal{S}$.\\
By induction, fix an isometry $\eta:\mathcal{A}^{\perp}\to\mathcal{B}^{\perp}$ satisfying $\eta|_{\mathcal{S}_0}=\sigma|_{\mathcal{S}_0}$. Then $(\sigma|_{\mathcal{A}})\odot\eta:\mathcal{V}\to\mathcal{W}$ extends $\sigma$.
\item \hlt{General Subspaces}: By Theorem \ref{thm:nonsingular-completions-hyperbolic-extensions}(iii), choose nonsingular extensions $\widehat{\mathcal{S}}\supseteq\mathcal{S}$ and $\widehat{\mathcal{T}}\supseteq\mathcal{T}$, isometry $\widehat{\sigma}:\widehat{\mathcal{S}}\to\widehat{\mathcal{T}}$ extending $\sigma$.\\
Item (iii) extends $\widehat{\sigma}$ to an isometry $\mathcal{V}\to\mathcal{W}$ and hence extends $\sigma$.
\end{enumerate}
\hlt{Structural Commentary}: Prepare elementary summands, induct through nonsingular subspaces, then complete arbitrary radicals.\\
\hlt{Key Insight}: Reflections handle orthogonal lines; hyperbolic planes handle symplectic pairs.
\end{proof}
\end{theorem}

\begin{corollary} \label{cor:witt-cancellation-metric-spaces} \hlt{Witt Cancellation}\\
Let $\mathcal{V}=\mathcal{S}\odot\mathcal{A}$, $\mathcal{W}=\mathcal{T}\odot\mathcal{B}$ be nonsingular metric vector spaces of the same kind, either symplectic or orthogonal with $\operatorname{char}(\mathbb{F})\neq2$.\\
If $\mathcal{V}\cong\mathcal{W}$ and $\mathcal{S}\cong\mathcal{T}$, then $\mathcal{A}\cong\mathcal{B}$.
\begin{proof}
Fix an isometry $\sigma:\mathcal{S}\to\mathcal{T}$. By Theorem \ref{thm:witt-extension-theorem-metric-spaces}, fix an isometry $\widehat{\sigma}:\mathcal{V}\to\mathcal{W}$ satisfying $\widehat{\sigma}|_{\mathcal{S}}=\sigma$.\\
Since $\mathcal{A}=\mathcal{S}^{\perp}$ and $\mathcal{B}=\mathcal{T}^{\perp}$, $\widehat{\sigma}(\mathcal{A})=\widehat{\sigma}(\mathcal{S}^{\perp})=(\widehat{\sigma}\mathcal{S})^{\perp}=\mathcal{T}^{\perp}=\mathcal{B}$. Hence $\widehat{\sigma}|_{\mathcal{A}}:\mathcal{A}\to\mathcal{B}$ is an isometry.\\
\hlt{Structural Commentary}: Extend the summand isometry and restrict the extension to its orthogonal complement.\\
\hlt{Key Insight}: Ambient extension converts cancellation into complement preservation.
\end{proof}
\end{corollary}

\begin{definition} \label{def:orthogonal-witt-index} \hlt{Witt Index}\\
Let $\mathcal{V}$ be nonsingular orthogonal with $\operatorname{char}(\mathbb{F})\neq2$. The \hlt{Witt index} is $w(\mathcal{V})\defeq\max\{\dim\mathcal{L}:\mathcal{L}\leq\mathcal{V},\ \mathcal{L}\subseteq\mathcal{L}^{\perp}\}$.
\end{definition}

\begin{theorem} \label{thm:orthogonal-witt-index-and-anisotropic-decomposition} \hlt{Witt Index and Anisotropic Decomposition}\\
Let $\mathcal{V}$ be orthogonal over $\mathbb{F}$ with $\operatorname{char}(\mathbb{F})\neq2$. Then:
\begin{enumerate}[label=(\roman*)]
\setlength{\itemsep}{0pt}
\item If $\mathcal{V}$ is nonsingular, all maximal totally isotropic subspaces have dimension $w(\mathcal{V})$.
\item If $\mathcal{V}$ is nonsingular, all maximal hyperbolic subspaces have dimension $2w(\mathcal{V})$.
\item If $R=\operatorname{rad}(\mathcal{V})$ and $\mathcal{U}$ is a nonsingular complement of $R$,\\
then $(\exists\,\mathcal{H},\mathcal{S}\leq\mathcal{U})\,\mathcal{H}$ is hyperbolic, $\mathcal{S}$ is anisotropic, $\mathcal{V}=\mathcal{S}\odot\mathcal{H}\odot R$, and $\dim\mathcal{H}=2w(\mathcal{U})$.
\end{enumerate}
\begin{proof}
\begin{enumerate}[label=(\roman*)]
\setlength{\itemsep}{0pt}
\item[$(i)$] Let $\mathcal{L}$ be maximal totally isotropic and choose $\mathcal{M}$ totally isotropic with $\dim\mathcal{M}=w(\mathcal{V})$.\\
Since $\dim\mathcal{L}\leq w(\mathcal{V})$, fix $\mathcal{M}_0\leq\mathcal{M}$ with $\dim\mathcal{M}_0=\dim\mathcal{L}$.\\
A vector-space isomorphism $\mathcal{L}\to\mathcal{M}_0$ is an isometry because both restricted forms vanish.\\
By Theorem \ref{thm:witt-extension-theorem-metric-spaces}, extend it to $\Phi\in\mathcal{O}(\mathcal{V})$. Then $\mathcal{L}\subseteq\Phi^{-1}\mathcal{M}$, so maximality proves $\mathcal{L}=\Phi^{-1}\mathcal{M}$ and $\dim\mathcal{L}=w(\mathcal{V})$.
\item[$(ii)$] A hyperbolic subspace with $k$ planes contains the $k$-dimensional totally isotropic span of one vector from each pair, so $k\leq w(\mathcal{V})$.\\
By Theorem \ref{thm:nonsingular-completions-hyperbolic-extensions}(i), a $w(\mathcal{V})$-dimensional totally isotropic subspace has a hyperbolic completion $\mathcal{H}_0$ of dimension $2w(\mathcal{V})$.\\
Suppose that a maximal hyperbolic $\mathcal{H}$ has dimension $2k<2w(\mathcal{V})$. Choose a $2k$-dimensional hyperbolic $\mathcal{H}_1\leq\mathcal{H}_0$.\\
By Theorem \ref{thm:witt-extension-theorem-metric-spaces}, extend an isometry $\mathcal{H}\to\mathcal{H}_1$ to $\Phi\in\mathcal{O}(\mathcal{V})$. \\
hen $\Phi^{-1}\mathcal{H}_0$ is a larger hyperbolic subspace containing $\mathcal{H}$, a contradiction.
\item[$(iii)$] By Theorem \ref{thm:radical-splitting-metric-space}, $\mathcal{V}=R\odot\mathcal{U}$ with $\mathcal{U}$ nonsingular. Choose a maximal hyperbolic $\mathcal{H}\leq\mathcal{U}$ and put $\mathcal{S}=\mathcal{H}^{\perp}\cap\mathcal{U}$.\\
By (ii) and Theorem \ref{thm:metric-orthogonal-complement-dimension}(ii), $\dim\mathcal{H}=2w(\mathcal{U})$ and $\mathcal{U}=\mathcal{H}\odot\mathcal{S}$. If $\mathcal{S}$ contained a nonzero isotropic vector, that vector would belong to a hyperbolic plane in $\mathcal{S}$ by Theorem \ref{thm:isotropic-vector-sits-in-hyperbolic-plane}, contradicting maximality of $\mathcal{H}$. Hence $\mathcal{S}$ is anisotropic and $\mathcal{V}=\mathcal{S}\odot\mathcal{H}\odot R$.
\end{enumerate}
\hlt{Structural Commentary}: Items (i)--(ii) identify the index geometrically; item (iii) splits the anisotropic residue.\\
\hlt{Key Insight}: The Witt index is the maximal number of hyperbolic planes.
\end{proof}
\end{theorem}
